\section{Arquitectura de implementación}
\label{sec:implementacion}

\subsection{Qué modelo está realmente implementado}
La formulación ejecutable del repositorio es \textbf{Benders/F4 con lazy cuts derivada de
F3}. En el informe se presentan F1, F2, F3 y F4 como linaje matemático, pero el solver C no expone
modelos monolíticos seleccionables para F1/F2/F3/F4. En particular, no se materializan variables
$x_{ij}$ de F1 ni variables $z_i^k$ de F2/F3 en el maestro C; se usan variables $y_j$ y
$\theta_i$ y se generan cortes de Benders de tipo F4 bajo demanda. El oráculo/prototipo Python y
la fuerza bruta se usan para validación en instancias pequeñas, no como benchmark monolítico
principal.

\subsection{Algoritmo de dos fases}
La implementación reproduce el \textbf{mecanismo central} del esquema de dos fases del paper
(Sección~3.5):
\begin{description}[leftmargin=1.5em]
  \item[Fase 1 — relajación lineal del maestro.] Se resuelve \eqref{eq:master} con $y\in[0,1]$,
        iterando: resolver el maestro LP, separar cortes violados (Alg.~1), agregarlos y aplicar
        una \textbf{heurística de redondeo} (abrir los $p$ sitios de mayor $\bar y_j$) para mejorar
        el $UB_1$. Termina cuando no hay más cortes violados: se obtiene $LB_1$ (valor de la
        relajación) y $UB_1$ (mejor solución entera redondeada).
  \item[Fase 2 — branch-and-Benders-cut.] Se imponen $y_j\in\{0,1\}$ y se resuelve con
        branch-and-cut. En cada \textbf{solución entera} que el solver encuentra, un
        \textbf{callback de lazy cuts} ejecuta el separador y agrega los cortes
        \eqref{eq:cut20} violados. Esto es, literalmente, branch-and-Benders-cut.
\end{description}
Como mejora, la Fase 2 \textbf{hereda el pool de cortes de la Fase 1} (warm-start): los cortes ya
encontrados se precargan como restricciones del MIP en lugar de partir vacíos.

\begin{quote}
\textbf{Precisión sobre el término ``warm-start''.} Aquí warm-start significa \emph{heredar los
cortes y cotas de la Fase 1} hacia la Fase 2 —transferencia de información entre fases, el
``warm start'' del Cap.~3—. \textbf{No} se refiere al \emph{preprocesamiento}: la construcción de
la matriz $S$ y de los radios ordenados $D^k_i$ (\texttt{sortsites\_build}) es una etapa previa,
común a ambas fases, independiente del warm-start. Tampoco es una heurística primal: solo
transfiere información dual/cortes al maestro MIP para que el árbol de branch-and-bound empiece
con una relajación más fuerte. Confundir estas piezas sería un error conceptual.
\end{quote}

\subsection{Mapa de las matemáticas a los módulos en C}
La implementación usa un \textbf{diseño modular con encapsulación por \texttt{struct}s e
interfaces}, inspirado en principios de orientación a objetos (C no es un lenguaje orientado a
objetos): un núcleo en C sin dependencias del solver y una capa delgada de solver (Gurobi) detrás
de una interfaz opaca (\texttt{src/solver.h}, \texttt{struct Solver}). El núcleo se valida contra
un oráculo en Python. La implementación es \emph{evidencia} de que la teoría funciona, no la
contribución del trabajo.

\begin{table}[h]\centering\small
\begin{tabular}{>{\raggedright}p{5.2cm} l l}
\toprule
\textbf{Objeto matemático} & \textbf{Archivo en C} & \textbf{Función / símbolo} \\
\midrule
Instancia, $d_{ij}$ (floor euclid./matriz) & \texttt{src/instance.c} & \texttt{instance\_load}, \texttt{instance\_dist} \\
Óptimo por fuerza bruta (oráculo) & \texttt{src/instance.c} & \texttt{instance\_brute\_force} \\
Matriz $S$, radios $D^k_i$ ordenados & \texttt{src/sortsites.c} & \texttt{sortsites\_build} $\to$ \texttt{SortSites} \\
$\kt_i$ — Algoritmo 2 & \texttt{src/separation.c} & \texttt{separation\_k\_tilde} \\
$\OPT(SP_i)$ \eqref{eq:opt18} + corte \eqref{eq:cut20} & \texttt{src/separation.c} & \texttt{separation\_client} \\
Algoritmo 1 (separar todos) & \texttt{src/separation.c} & \texttt{separation\_all} (\texttt{cut\_sink}) \\
Heurística de redondeo & \texttt{src/heuristic.c} & \texttt{rounding\_heuristic} \\
Fase 1 (maestro LP + loop) & \texttt{src/phase1.c} & \texttt{phase1\_run} \\
Fase 2 (B\&Benders-cut + callback) & \texttt{src/phase2.c} & \texttt{phase2\_run}, \texttt{p2\_lazy\_fn} \\
Pool de cortes (warm-start) & \texttt{src/cutpool.c} & \texttt{cutpool\_add} \\
Capa de solver (abstracción) & \texttt{src/solver\_gurobi.c} & \texttt{solver\_*}, \texttt{GRBcblazy} \\
CLI, logging, CSV & \texttt{src/main.c}, \texttt{src/logging.c} & \texttt{main}, \texttt{csv\_append} \\
\midrule
Oráculo de correctitud & \texttt{prototype/pmp\_benders.py} & \texttt{separate\_client}, \texttt{phase2} \\
Verificación del separador & \texttt{tests/test\_separation\_toy.c} & \texttt{make test} \\
Cross-check C\,$\equiv$\,Python & \texttt{scripts/verify\_cuts.py} & — \\
\bottomrule
\end{tabular}
\caption{Trazabilidad de cada objeto del modelo a su implementación.}
\end{table}

\subsection{Detalles de modelamiento}
\begin{itemize}[noitemsep]
  \item \textbf{Variables del maestro:} $y_j$ con índice $j=0,\dots,M-1$; $\theta_i$ con índice
        $M+i$. Objetivo $\min\sum_i\theta_i$; única restricción estructural $\sum_j y_j=p$. Los
        cortes \eqref{eq:cut20} se agregan como filas $\theta_i + \sum_t \text{coef}_t\,y_{j_t}\ge
        \text{rhs}$.
  \item \textbf{Callback:} en \texttt{GRB\_CB\_MIPSOL} se lee la solución entera candidata, se
        corre \texttt{separation\_all} y cada corte violado se inyecta con \texttt{GRBcblazy}
        (parámetro \texttt{LazyConstraints=1}).
  \item \textbf{Distancia:} euclidiana \emph{truncada} (floor) para coordenadas (convención del
        paper, validada numéricamente en \texttt{rl1304}); lectura directa para instancias con
        matriz (incluida la asimétrica RW). Distancias como enteros largos; objetivo como
        \texttt{double}.
  \item \textbf{Exactitud:} \texttt{MIPGap}$=10^{-10}$ para probar optimalidad exacta.
\end{itemize}

\subsection{Capa de solver y portabilidad}
Toda la interacción con Gurobi pasa por una interfaz delgada (\texttt{src/solver.h}) que aísla el
backend: agregar variables/restricciones, optimizar, leer solución, contar nodos y manejar el
callback de lazy cuts. Esto deja el núcleo (parsing, distancias, $S$, separación, fases)
independiente del solver y permitiría, en principio, intercambiar el backend (CPLEX/SCIP) sin
tocar la lógica de Benders. En esta entrega solo se implementó y probó el backend Gurobi; CPLEX y
SCIP quedan como extensiones futuras.
